\scanpage{10}\mspage{4}
\[
A=\{\varepsilon x,\varepsilon'y,-(\varepsilon x+\varepsilon'y)\},\qquad \varepsilon,\varepsilon'\in\{\pm1_{\deleted{\unsure{\Lambda}}}\}.
\]
\deleted{donc l'on est ramené à \unsure{choi.|choisir}} \inserted{et l'ens. des $f$ cherchés} s'identifie au quotient par l'involution $(\varepsilon,\varepsilon')\mapsto(-\varepsilon,-\varepsilon')$, de l'ens des $(\varepsilon,\varepsilon')\in\{\pm1\}_{\deleted{\unsure{\Lambda}}}^2$. On a donc exactement deux tels $f$. En résumé :

\uline{Proposition}\quad Soient $\Lambda$ un anneau tel que $2.1_\Lambda\ne0$, et $M$ un $\Lambda$-module libre de rang $2$. Soient \deleted{$M^*=\{x\in M\mid M/\Lambda x$} \deleted{$\widetilde S_M=\{x\in M$}
% slanted
\margin{on a donc modifié légèrement la définition (15)}
\[
\left\{\begin{aligned}
\widetilde S_M&=\{x\in M\mid\exists y\in M\ \text{tel que }\unsure{\{|(}x,y)\in\mathrm{Rep}(M)\}\\
\deleted{S_M}&\deleted{=\widetilde S_M/\{\pm\mathrm{id}_M\}}\\
\widetilde F_M&=\left\{\{x,y,z\}\in\unsure{\mathfrak P|\mathscr P}_3(M)\ \middle|\
\begin{gathered}x+y+z=0\\\unsure{\{|(}x\unsure{+|,}y\unsure{\}|)}\in\mathrm{Rep}(M)\end{gathered}\right\},\\[-2pt]
&\hspace{3em}\text{(i.e. }(y,z)\in\mathrm{Rep}(M),\\
&\hspace{3em}\text{i.e. }(z,x)\in\mathrm{Rep}(M).
\end{aligned}\right.\tag*{(20)}
\]
\[
\left\{\begin{aligned}
S_M&=\widetilde S_M/\{\pm1\}\qquad p_M:\ \widetilde S_M\to S_M\quad\text{appl. can.}\\
F_M&=\left\{f\in\mathfrak P_{\unsure{3|2}}(S_M)\ \middle|\
\begin{gathered}\exists\widetilde f\in\widetilde F_M,\ \text{tel que}\\ f=p(\widetilde f)\end{gathered}\right\},\\
A_M&=\{a\in\mathfrak P_2(S_M)\mid\exists f\in F_M,\ a\subset f\}
\end{aligned}\right.\tag*{(21)}
\]
Alors $\widetilde F_M$ est formé de parties à trois éléments de $\widetilde S_M$, de réunion $\widetilde S_M$, $F_M$ est formé de parties à trois éléments de $S_M$, de réunion $S_M$, et l'application \inserted{surj.} $\widetilde F_M\to F_M$, $\widetilde f\to p(\widetilde f)$,
% archive sheet number: 9
